% Tapis de Quillen, with Quillen's offprint and preprint (cote 162-5, typed part): archive cote 162-5, site page 36 (PDF page 37).
% typescript. Draft compared: none.
% AI-transcribed by Claude Opus 5.5 (claude-opus-5-5, adaptive thinking, effort high), Message Batches API, batch t162_full (msgbatch_01UAh3zyWPd454AeDkhnmQ33), 2026-11-13; not yet checked by eye.
\scanpage{162-5/36}
éventuellement muni de la sous-catégorie pleine de tous les systêmes locau
sur X, et du foncteur cohomologique sur $D^b_{lc}(X)$ à valeurs dans \unsure{ladite}
catégorie, et bien sûr du produit tensoriel (mais alors on sort de $D^b$
pour entrer dans $D^-$ , redactor demerdetur).

II n-catégories, catégories n-uples, et (Gr)-catégories.

Appelons bi-catégorie une catégorie dans(Cat), i.e. un triple
$(C^{\deleted{o}}_{\hand{0}},C^{\deleted{1}}_{\hand{1}},$ \xout{\unsure{xxoù}}$p_1,p_2,\hand{\pi}$ ), où $C^o$ est une catégorie, $p_1,p_2$: $C^{\deleted{1}}_{\hand{1}}\hand{\rightrightarrows}C^{\deleted{o}}_{\hand{0}}$ deux
foncteurs, $\hand{\pi}:(C^{\deleted{1}}_{\hand{1}},p_1)x_{C^o}(C^{\deleted{1}}_{\hand{1}},p_2)\hand{\longrightarrow}C^{\deleted{1}}_{\hand{1}}$ un foncteur, avec les axiomes
habituels. Interprétant (Cat) comme la sous-catégorie pleine de (Ssimpl)
des objets qui .... , et de même pour les catégories à valeurs dans une
catégorie telle (Cat), on trouve que la catégorie \inserted[typed]{(Bicat)} des (Cat)-catégories
est équivalente à la sous-catégorie pleine de la catégorie des ensembles
bi-semisimpliciaux

\diagram{
\[
\begin{array}{llll}
X^{\deleted{00}}_{\hand{00}} & X^{\deleted{01}}_{\hand{01}} & \ldots & X^{\deleted{0,n}}_{\hand{0n}}\\
X^{\deleted{10}}_{\hand{10}} & X^{\deleted{11}}_{\hand{11}} & \ldots & X^{\deleted{1,n}}_{\hand{1n}}\\
\multicolumn{4}{l}{\ldots\ldots\ldots}
\end{array}
\]
}

formée des objets qui ligne par ligne et colonne par colonne satisfoht la
condition d'exctitude familière (transformer sommes amalgamées en produits
fibrés). De même la catégorie \xout{des} (Tricat) des (Bicat)-catégories
s'interprète en t ermes d'objets tri-\inserted[typed]{semi}simpliciaux, etc.

Exemple: soit C une catégorie, posons $C^{\deleted{o}}_{\hand{0}}=C$, $C^{\deleted{1}}_{\hand{1}}=\underline{Fl}(C)$, avec structures
$p_1,p_2$, évidentes (d'où une (Cat)-catétorie dont les composantes semi-sim-%
pliciales supérieures sont les catégories $\underline{Fl}^n(C)$ des n-flêches de C ...).
L'ensemble bi-semisimplicial correspondant a pour composantes les ensembles
Hom($\hand{\Delta}^n$x$\Delta^m$,C). Par exemple, les objets de $X^{\deleted{11}}_{\hand{11}}$ sont les carrés commutatifs
de flêches de C, qu'on peut interpréter(de deux manières différentes) comme

\mspage{\hand{36}}
